\documentclass[12pt,a4paper]{report}
\usepackage[utf8]{inputenc}
\usepackage[T1]{fontenc}
\usepackage[french]{babel}
\usepackage{geometry}
\usepackage{xcolor}
\usepackage{tabularx}
\usepackage{booktabs}
\usepackage{longtable}
\usepackage{enumitem}
\usepackage{fancyhdr}
\usepackage{titlesec}
\usepackage{hyperref}
\usepackage{listings}
\usepackage{array}
\usepackage{multirow}
\usepackage{tcolorbox}
\usepackage{mdframed}
\usepackage{tikz}
\usepackage{pgfplots}
\pgfplotsset{compat=1.18}

\geometry{top=2.5cm, bottom=2.5cm, left=3cm, right=2.5cm}

%----------- Couleurs -----------
\definecolor{primary}{RGB}{30,41,59}
\definecolor{accent}{RGB}{59,130,246}
\definecolor{success}{RGB}{16,185,129}
\definecolor{warning}{RGB}{245,158,11}
\definecolor{danger}{RGB}{239,68,68}
\definecolor{lightgray}{RGB}{248,250,252}
\definecolor{bordergray}{RGB}{226,232,240}
\definecolor{codebg}{RGB}{15,23,42}
\definecolor{codetext}{RGB}{226,232,240}
\definecolor{codekw}{RGB}{147,197,253}
\definecolor{codestr}{RGB}{134,239,172}
\definecolor{codecomment}{RGB}{148,163,184}

%----------- Hyperlinks -----------
\hypersetup{
    colorlinks=true,
    linkcolor=accent,
    urlcolor=accent,
    pdftitle={Document d'Architecture Technique - GRM},
    pdfauthor={Équipe GRM}
}

%----------- En-tête / Pied de page -----------
\pagestyle{fancy}
\fancyhf{}
\fancyhead[L]{\small\color{primary}\textbf{GRM} — Document d'Architecture Technique}
\fancyhead[R]{\small\color{primary}\leftmark}
\fancyfoot[L]{\small\color{gray}Confidentiel — Usage interne}
\fancyfoot[C]{\thepage}
\fancyfoot[R]{\small\color{gray}\today}
\renewcommand{\headrulewidth}{0.5pt}
\renewcommand{\footrulewidth}{0.3pt}

%----------- Titres de chapitres -----------
\titleformat{\chapter}[block]
  {\normalfont\Large\bfseries\color{primary}}
  {\thechapter.}{0.8em}{}[\vspace{-4pt}\color{accent}\rule{\textwidth}{2pt}\vspace{6pt}]
\titlespacing*{\chapter}{0pt}{0pt}{16pt}
\titleformat{\section}{\normalfont\large\bfseries\color{primary}}{\thesection}{0.8em}{}
\titleformat{\subsection}{\normalfont\normalsize\bfseries\color{primary}}{\thesubsection}{0.8em}{}
\titleformat{\subsubsection}{\normalfont\small\bfseries\color{primary}}{\thesubsubsection}{0.8em}{}

%----------- Listings -----------
\lstdefinestyle{sql}{
    language=SQL,
    backgroundcolor=\color{codebg},
    basicstyle=\ttfamily\small\color{codetext},
    keywordstyle=\color{codekw}\bfseries,
    stringstyle=\color{codestr},
    commentstyle=\color{codecomment}\itshape,
    numbers=left,
    numberstyle=\tiny\color{gray},
    breaklines=true,
    frame=none,
    tabsize=2,
    showstringspaces=false,
    captionpos=b,
    rulecolor=\color{bordergray}
}

\lstdefinestyle{js}{
    language=JavaScript,
    backgroundcolor=\color{codebg},
    basicstyle=\ttfamily\small\color{codetext},
    keywordstyle=\color{codekw}\bfseries,
    stringstyle=\color{codestr},
    commentstyle=\color{codecomment}\itshape,
    numbers=left,
    numberstyle=\tiny\color{gray},
    breaklines=true,
    frame=none,
    tabsize=2,
    showstringspaces=false,
    captionpos=b
}

\lstdefinestyle{bash}{
    language=bash,
    backgroundcolor=\color{codebg},
    basicstyle=\ttfamily\small\color{codetext},
    keywordstyle=\color{codekw}\bfseries,
    stringstyle=\color{codestr},
    commentstyle=\color{codecomment}\itshape,
    numbers=none,
    breaklines=true,
    frame=none,
    tabsize=2,
    showstringspaces=false,
    captionpos=b
}

%----------- Boîtes personnalisées -----------
\tcbset{
  infobox/.style={
    colback=lightgray,
    colframe=accent,
    fonttitle=\bfseries\small\color{primary},
    boxrule=1pt,
    arc=4pt,
    left=8pt, right=8pt, top=6pt, bottom=6pt
  },
  warnbox/.style={
    colback={RGB:orange,5;white,95},
    colframe=warning,
    fonttitle=\bfseries\small\color{primary},
    boxrule=1pt,
    arc=4pt,
    left=8pt, right=8pt, top=6pt, bottom=6pt
  }
}

\newcolumntype{L}[1]{>{\raggedright\arraybackslash}p{#1}}
\newcolumntype{C}[1]{>{\centering\arraybackslash}p{#1}}

%==========================================================
\begin{document}
%==========================================================

%----------- PAGE DE GARDE -----------
\begin{titlepage}
\pagecolor{primary}
\color{white}
\vspace*{2cm}
\begin{center}
    {\fontsize{38}{46}\selectfont\bfseries GRM}\\[0.4cm]
    {\Large\bfseries Gestion des Ressources Matérielles}\\[2.5cm]

    \colorbox{accent}{\parbox{0.85\textwidth}{\centering\vspace{0.4cm}
    {\huge\bfseries Document d'Architecture Technique}\\[0.2cm]
    {\large Conception, Technologies \& Déploiement}
    \vspace{0.4cm}}}

    \vspace{2.5cm}

    \begin{tabular}{ll}
        \textbf{Projet} & Gestion des Ressources Matérielles \\[4pt]
        \textbf{Établissement} & Ynov Campus \\[4pt]
        \textbf{Filière} & M1 Informatique — Architecture Logicielle \\[4pt]
        \textbf{Responsable pédagogique} & Pr. Driss \\[4pt]
        \textbf{Version} & 1.0 \\[4pt]
        \textbf{Date} & \today \\[4pt]
        \textbf{Statut} & \textcolor{success}{Approuvé} \\
    \end{tabular}

    \vspace{3cm}
    {\small\color{gray} Ce document est confidentiel et destiné à un usage interne uniquement.}
\end{center}
\end{titlepage}
\pagecolor{white}\color{black}

%----------- HISTORIQUE DES RÉVISIONS -----------
\chapter*{Historique des Révisions}
\addcontentsline{toc}{chapter}{Historique des Révisions}
\begin{tabularx}{\textwidth}{|C{1.8cm}|C{2.5cm}|L{4.5cm}|L{4.5cm}|}
\hline
\rowcolor{primary}\textcolor{white}{\textbf{Version}} & \textcolor{white}{\textbf{Date}} & \textcolor{white}{\textbf{Description}} & \textcolor{white}{\textbf{Auteur}} \\
\hline
0.1 & 01/04/2025 & Rédaction de la vue d'ensemble & Équipe GRM \\
\hline
0.5 & 20/04/2025 & Ajout des diagrammes et choix technologiques & Équipe GRM \\
\hline
0.9 & 05/05/2025 & Ajout des diagrammes de séquence et déploiement & Équipe GRM \\
\hline
1.0 & \today & Version finale validée & Équipe GRM \\
\hline
\end{tabularx}

%----------- TABLE DES MATIÈRES -----------
\tableofcontents
\newpage

%==========================================================
\chapter{Vue d'Ensemble de l'Architecture}
%==========================================================

\section{Introduction}

Ce document décrit l'architecture technique complète du système GRM. Il couvre les décisions architecturales, les choix technologiques, la conception des composants, les flux de communication, la sécurité et les stratégies de déploiement.

L'architecture du système GRM répond aux contraintes suivantes :
\begin{itemize}[leftmargin=1.5cm]
    \item \textbf{Séparation des préoccupations} : chaque couche a une responsabilité claire et délimitée
    \item \textbf{Sécurité par conception} : authentification et autorisation intégrées dès la conception
    \item \textbf{Maintenabilité} : code modulaire, configuration externalisée, déploiement reproductible
    \item \textbf{Extensibilité} : ajout de nouveaux modules sans remise en cause de l'architecture existante
\end{itemize}

\section{Pattern Architectural : 3-Tiers}

Le système adopte le pattern \textbf{3-tiers} (three-tier architecture), qui sépare le système en trois couches logiques indépendantes :

\begin{tcolorbox}[infobox, title=Définition — Architecture 3-Tiers]
L'architecture 3-tiers est un modèle de conception logicielle qui divise une application en trois couches distinctes :
\begin{enumerate}
    \item \textbf{Couche Présentation (Tier 1)} : Interface utilisateur — gère l'affichage et les interactions
    \item \textbf{Couche Logique Métier (Tier 2)} : Traitement — applique les règles métier et orchestre les données
    \item \textbf{Couche Données (Tier 3)} : Persistance — stocke et récupère les données de manière structurée
\end{enumerate}
\end{tcolorbox}

\section{Diagramme d'Architecture Global}

\begin{center}
\begin{verbatim}
  ┌─────────────────────────────────────────────────────────┐
  │                   TIER 1 — PRÉSENTATION                 │
  │                                                         │
  │   ┌──────────────────────────────────────────────┐      │
  │   │       React SPA — Application Web            │      │
  │   │  Port 3000 (dev) / Port 80 (prod via Nginx)  │      │
  │   │                                              │      │
  │   │  AuthContext │ Pages │ Composants │ Services │      │
  │   └──────────────────────────────────────────────┘      │
  └────────────────────────┬────────────────────────────────┘
                           │ HTTP/HTTPS — REST API (JSON)
                           │ En-tête: Authorization: Bearer <JWT>
  ┌────────────────────────▼────────────────────────────────┐
  │                   TIER 2 — LOGIQUE MÉTIER               │
  │                                                         │
  │   ┌──────────────────────────────────────────────┐      │
  │   │    Express.js API Server — Node.js 20        │      │
  │   │              Port 5000                       │      │
  │   │                                              │      │
  │   │  Middleware JWT │ Routes │ Controllers       │      │
  │   └──────────────────────────────────────────────┘      │
  └────────────────────────┬────────────────────────────────┘
                           │ SQL — Requêtes paramétrées
                           │ Protocole TCP/5432 (pg pool)
  ┌────────────────────────▼────────────────────────────────┐
  │                   TIER 3 — DONNÉES                      │
  │                                                         │
  │   ┌──────────────────────────────────────────────┐      │
  │   │         PostgreSQL 15 — Base de données      │      │
  │   │              Port 5432                       │      │
  │   │                                              │      │
  │   │  11 tables │ UUID │ JSONB │ Transactions     │      │
  │   └──────────────────────────────────────────────┘      │
  └─────────────────────────────────────────────────────────┘
\end{verbatim}
\end{center}

\section{Justification des Choix Architecturaux}

\subsection{Pourquoi une Architecture 3-Tiers ?}

\begin{longtable}{|L{3.5cm}|L{10cm}|}
\caption{Avantages de l'architecture 3-tiers pour GRM} \\
\hline
\rowcolor{primary}\textcolor{white}{\textbf{Aspect}} & \textcolor{white}{\textbf{Justification}} \\
\hline
\endfirsthead
\hline
\textbf{Aspect} & \textbf{Justification} \\
\hline
\endhead
Séparation claire & Le frontend peut évoluer (changement de framework) sans toucher au backend, et vice-versa. \\
\hline
Réutilisation & L'API REST peut être consommée par d'autres clients (application mobile, scripts d'administration) sans modification. \\
\hline
Sécurité & La logique de contrôle d'accès est centralisée dans le backend, impossible à contourner côté client. \\
\hline
Testabilité & Chaque couche peut être testée de manière indépendante (tests unitaires, tests d'intégration). \\
\hline
Scalabilité & Chaque tier peut être scalé indépendamment selon la charge (ex: plusieurs instances backend avec une seule DB). \\
\hline
\end{longtable}

\subsection{Alternatives Considérées}

\begin{longtable}{|L{3cm}|L{5.5cm}|L{5.5cm}|}
\caption{Comparaison des alternatives architecturales} \\
\hline
\rowcolor{primary}\textcolor{white}{\textbf{Alternative}} & \textcolor{white}{\textbf{Avantages}} & \textcolor{white}{\textbf{Inconvénients / Raison de rejet}} \\
\hline
\endfirsthead
\hline
\textbf{Alternative} & \textbf{Avantages} & \textbf{Inconvénients / Raison de rejet} \\
\hline
\endhead
Architecture Monolithique & Simple à déployer, faible overhead réseau & Couplage fort, difficultés de maintenance, scalabilité limitée \\
\hline
Microservices & Isolation parfaite, scalabilité fine & Complexité opérationnelle excessive pour la taille du projet \\
\hline
MVC Serveur (SSR) & Rendu côté serveur, SEO natif & Moins adapté à une interface riche et temps réel; rechargement complet de page \\
\hline
GraphQL & Flexibilité des requêtes, typage & Surcharge de configuration, courbe d'apprentissage plus élevée \\
\hline
\end{longtable}

%==========================================================
\chapter{Choix Technologiques}
%==========================================================

\section{Stack Technologique Complète}

\begin{longtable}{|L{2.5cm}|L{2.5cm}|L{3.5cm}|L{5cm}|}
\caption{Stack technologique GRM} \\
\hline
\rowcolor{primary}\textcolor{white}{\textbf{Couche}} & \textcolor{white}{\textbf{Technologie}} & \textcolor{white}{\textbf{Version}} & \textcolor{white}{\textbf{Justification}} \\
\hline
\endfirsthead
\hline
\textbf{Couche} & \textbf{Technologie} & \textbf{Version} & \textbf{Justification} \\
\hline
\endhead
Frontend UI & React & 18.2 & SPA déclarative, composants réutilisables, large communauté \\
\hline
Routing SPA & React Router & v6 & Navigation client-side sans rechargement, gestion de routes imbriquées \\
\hline
HTTP Client & Axios & 1.6.x & Intercepteurs pour gestion centralisée du JWT, meilleure ergonomie que fetch natif \\
\hline
Build Tool & Create React App & 5.0 & Configuration webpack zero-config, intégration ESLint \\
\hline
Backend & Node.js & 20 LTS & Non-bloquant (event loop), JSON natif, même langage côté client/serveur \\
\hline
Framework API & Express.js & 4.18 & Minimaliste, flexible, middleware system mature \\
\hline
Auth & jsonwebtoken & 9.x & Standard JWT RFC 7519, signature HMAC-SHA256 \\
\hline
Hachage & bcrypt & 5.1 & Standard de l'industrie pour hachage sécurisé des mots de passe \\
\hline
Base de données & PostgreSQL & 15 & ACID, JSONB natif, UUID, transactions imbriquées, robustesse prouvée \\
\hline
Driver DB & pg (node-postgres) & 8.11 & Pool de connexions intégré, requêtes paramétrées, support des transactions \\
\hline
Proxy / Reverse & Nginx & 1.25 (Alpine) & Sert le build statique React et proxifie les appels API vers le backend \\
\hline
Conteneurisation & Docker + Compose & 24.x / 2.x & Reproductibilité des environnements, isolation, déploiement unifié \\
\hline
\end{longtable}

\section{Justification Détaillée des Choix}

\subsection{Node.js + Express vs Spring Boot}

\begin{tcolorbox}[infobox, title=Pourquoi Node.js plutôt que Spring Boot ?]
\begin{itemize}
    \item \textbf{Homogénéité du langage} : JavaScript côté client et serveur réduit le contexte de commutation pour les développeurs
    \item \textbf{Légèreté} : Express est minimaliste (~50 KB) contre Spring Boot (~15 MB), réduisant l'empreinte mémoire du conteneur
    \item \textbf{Écosystème npm} : plus de 2 millions de packages disponibles
    \item \textbf{Performance I/O} : le modèle non-bloquant est idéal pour une API REST majoritairement I/O-bound (lectures DB)
    \item \textbf{Démarrage rapide} : le serveur Node.js démarre en < 500 ms, contre 10-30 secondes pour Spring Boot
\end{itemize}
\textit{Spring Boot aurait été préféré pour un système avec logique de calcul intensive (CPU-bound) ou un équipe Java expérimentée.}
\end{tcolorbox}

\subsection{PostgreSQL vs MySQL vs MongoDB}

\begin{longtable}{|L{3cm}|L{3.5cm}|L{3.5cm}|L{3.5cm}|}
\caption{Comparaison des SGBD} \\
\hline
\rowcolor{primary}\textcolor{white}{\textbf{Critère}} & \textcolor{white}{\textbf{PostgreSQL}} & \textcolor{white}{\textbf{MySQL}} & \textcolor{white}{\textbf{MongoDB}} \\
\hline
\endfirsthead
\hline
\textbf{Critère} & \textbf{PostgreSQL} & \textbf{MySQL} & \textbf{MongoDB} \\
\hline
\endhead
Transactions ACID & Complet & Partiel (InnoDB) & Limité \\
\hline
JSONB natif & Oui (indexable) & Non & Oui (natif) \\
\hline
UUID natif & Oui & Non & Oui \\
\hline
Contraintes FK & Robustes & Robustes & Aucune \\
\hline
Typage fort & Oui & Oui & Non \\
\hline
Intégrité données & Haute & Haute & Faible \\
\hline
Choix final & \textbf{Retenu} & Alternative & Non adapté \\
\hline
\end{longtable}

PostgreSQL a été retenu car il offre à la fois la rigueur des contraintes relationnelles (clés étrangères, contraintes CHECK, transactions ACID) et la flexibilité du JSONB pour les spécifications techniques variables des ressources.

\subsection{React vs Angular}

React a été préféré à Angular pour les raisons suivantes :
\begin{itemize}[leftmargin=1.5cm]
    \item \textbf{Courbe d'apprentissage} plus douce (bibliothèque vs framework complet)
    \item \textbf{Flexibilité} : React impose moins de conventions, laissant plus de liberté architecturale
    \item \textbf{Performance} : Virtual DOM optimisé, rendu différentiel efficace
    \item \textbf{Écosystème} : taille de la communauté et richesse des bibliothèques tierces
\end{itemize}

%==========================================================
\chapter{Architecture du Backend}
%==========================================================

\section{Organisation du Code}

L'API backend suit le pattern \textbf{MVC allégé} (Model-View-Controller), sans couche View (remplacée par des réponses JSON).

\begin{center}
\begin{verbatim}
backend/
├── server.js                  ← Point d'entrée — configuration Express,
│                                montage des routes, démarrage serveur
├── .env                       ← Variables de configuration (non versionné)
├── package.json
├── Dockerfile
│
├── database/
│   └── schema.sql             ← DDL complet + données initiales
│
└── src/
    ├── config/
    │   └── database.js        ← Pool pg (singleton)
    │
    ├── middleware/
    │   └── auth.js            ← authenticate() + authorize(...roles)
    │
    ├── controllers/           ← Logique métier (1 fichier par domaine)
    │   ├── authController.js
    │   ├── tenderController.js
    │   ├── offerController.js
    │   ├── resourceController.js
    │   ├── supplierController.js
    │   ├── breakdownController.js
    │   └── departmentController.js
    │
    └── routes/                ← Définition des routes REST + application middleware
        ├── auth.js
        ├── tenders.js
        ├── offers.js
        ├── resources.js
        ├── suppliers.js
        ├── breakdowns.js
        └── departments.js
\end{verbatim}
\end{center}

\section{Flux d'une Requête HTTP}

Toute requête entrante suit le flux suivant :

\begin{center}
\begin{verbatim}
Client HTTP
    │
    ▼
[1] server.js
    ├── CORS Middleware         → Vérifie l'origine de la requête
    ├── express.json()         → Parse le corps JSON
    └── Montage des routes     → /api/auth, /api/tenders, etc.
    │
    ▼
[2] Router (ex: routes/tenders.js)
    ├── authenticate           → Vérifie et décode le JWT
    ├── authorize('responsable')  → Contrôle le rôle
    └── Route Handler          → Appelle le controller
    │
    ▼
[3] Controller (ex: tenderController.js)
    ├── Extraction paramètres  → req.body, req.params, req.user
    ├── Validation métier      → Règles de gestion
    ├── Requête SQL (pool.query)  → Accès PostgreSQL
    └── Réponse JSON           → res.json() ou res.status(xxx).json()
    │
    ▼
[4] PostgreSQL
    ├── Pool de connexions     → Gestion automatique
    ├── Transaction (si besoin)  → BEGIN / COMMIT / ROLLBACK
    └── Retourne les résultats
\end{verbatim}
\end{center}

\section{Conception de l'API REST}

\subsection{Conventions REST Respectées}

\begin{longtable}{|L{2.5cm}|L{2cm}|L{4cm}|L{5cm}|}
\caption{Endpoints de l'API GRM} \\
\hline
\rowcolor{primary}\textcolor{white}{\textbf{Méthode}} & \textcolor{white}{\textbf{Code}} & \textcolor{white}{\textbf{Route}} & \textcolor{white}{\textbf{Action}} \\
\hline
\endfirsthead
\hline
\textbf{Méthode} & \textbf{Code} & \textbf{Route} & \textbf{Action} \\
\hline
\endhead
POST & 200 & /api/auth/login & Authentification \\
\hline
POST & 201 & /api/auth/register-supplier & Inscription fournisseur \\
\hline
GET & 200 & /api/auth/profile & Profil connecté \\
\hline
GET & 200 & /api/tenders & Liste des appels d'offre \\
\hline
POST & 201 & /api/tenders & Créer un appel d'offre \\
\hline
GET & 200 & /api/tenders/:id & Détail d'un appel d'offre \\
\hline
PUT & 200 & /api/tenders/:id & Modifier un appel d'offre \\
\hline
PATCH & 200 & /api/tenders/:id/close & Fermer un appel d'offre \\
\hline
GET & 200 & /api/offers/tender/:id & Offres d'un AO \\
\hline
POST & 201 & /api/offers/tender/:id & Soumettre une offre \\
\hline
PATCH & 200 & /api/offers/:id/accept & Accepter une offre \\
\hline
PATCH & 200 & /api/offers/:id/reject & Rejeter une offre \\
\hline
GET & 200 & /api/resources & Liste des ressources \\
\hline
POST & 201 & /api/resources & Créer une ressource \\
\hline
GET & 200 & /api/resources/:id & Détail d'une ressource \\
\hline
PUT & 200 & /api/resources/:id & Modifier une ressource \\
\hline
DELETE & 200 & /api/resources/:id & Supprimer une ressource \\
\hline
POST & 201 & /api/resources/:id/assign & Affecter une ressource \\
\hline
GET & 200 & /api/suppliers & Liste des fournisseurs \\
\hline
PATCH & 200 & /api/suppliers/:id/blacklist & Blacklister \\
\hline
PATCH & 200 & /api/suppliers/:id/unblacklist & Retirer du blacklist \\
\hline
GET & 200 & /api/breakdowns & Liste des pannes \\
\hline
POST & 201 & /api/breakdowns & Signaler une panne \\
\hline
GET & 200 & /api/breakdowns/:id & Détail d'une panne \\
\hline
POST & 201 & /api/breakdowns/:id/maintenance-report & Constat technicien \\
\hline
PATCH & 200 & /api/breakdowns/:id/resolve & Résoudre une panne \\
\hline
PATCH & 200 & /api/breakdowns/:id/return-to-supplier & Retour fournisseur \\
\hline
GET & 200 & /api/departments & Liste des départements \\
\hline
POST & 201 & /api/departments & Créer un département \\
\hline
GET & 200 & /api/departments/:id/members & Membres d'un département \\
\hline
GET & 200 & /api/health & Santé du serveur \\
\hline
\end{longtable}

\subsection{Format des Réponses}

\textbf{Succès (200/201) :}
\begin{lstlisting}[style=js, caption={Format réponse succès}]
// Ressource unique
{ "id": "uuid", "name": "...", "status": "available", ... }

// Liste
[ { "id": "uuid", ... }, { "id": "uuid", ... } ]

// Message de confirmation
{ "message": "Resource deleted" }
\end{lstlisting}

\textbf{Erreur (4xx/5xx) :}
\begin{lstlisting}[style=js, caption={Format réponse erreur}]
// 401 - Non authentifié
{ "error": "No token provided" }
{ "error": "Invalid token" }

// 403 - Non autorisé
{ "error": "Access denied" }
{ "error": "Supplier is blacklisted" }

// 404 - Non trouvé
{ "error": "Resource not found" }

// 500 - Erreur serveur
{ "error": "duplicate key value violates unique constraint..." }
\end{lstlisting}

\section{Gestion des Transactions}

Les opérations critiques qui modifient plusieurs tables simultanément utilisent des transactions SQL explicites pour garantir l'atomicité.

\begin{lstlisting}[style=js, caption={Transaction — Acceptation d'une offre (offerController.js)}]
const accept = async (req, res) => {
    const client = await pool.connect();
    try {
        await client.query('BEGIN');

        // 1. Marquer l'offre selectionnee comme acceptee
        const { rows } = await client.query(
            `UPDATE offers SET status='accepted'
             WHERE id = $1 RETURNING *`,
            [req.params.id]
        );

        // 2. Rejeter automatiquement toutes les autres offres
        await client.query(
            `UPDATE offers SET status='rejected'
             WHERE tender_id = $1 AND id != $2`,
            [rows[0].tender_id, req.params.id]
        );

        // 3. Marquer l'appel d'offre comme attribue
        await client.query(
            `UPDATE tenders SET status='awarded'
             WHERE id = $1`,
            [rows[0].tender_id]
        );

        await client.query('COMMIT');
        res.json(rows[0]);
    } catch (err) {
        await client.query('ROLLBACK');
        res.status(500).json({ error: err.message });
    } finally {
        client.release(); // Retour au pool obligatoire
    }
};
\end{lstlisting}

%==========================================================
\chapter{Architecture du Frontend}
%==========================================================

\section{Organisation du Code React}

\begin{center}
\begin{verbatim}
frontend/src/
├── index.js                   ← Point d'entrée — ReactDOM.createRoot()
├── App.js                     ← Router principal (BrowserRouter + Routes)
│
├── context/
│   └── AuthContext.js         ← État global d'authentification (React Context)
│                                Expose: user, login(), logout(), loading
│
├── services/
│   └── api.js                 ← Instance Axios configurée
│                                Intercepteur: redirection sur 401
│
├── components/
│   └── Layout/
│       └── Layout.js          ← Shell de l'application:
│                                Sidebar (collapsible) + Header + <Outlet>
│
└── pages/                     ← Pages par domaine fonctionnel
    ├── Login.js               ← Formulaire d'authentification
    ├── Dashboard.js           ← Vue synthétique avec statistiques
    ├── Tenders/
    │   ├── TenderList.js      ← Liste des appels d'offre
    │   ├── TenderForm.js      ← Formulaire de création
    │   └── TenderDetail.js    ← Détail + gestion des offres
    ├── Resources/
    │   ├── ResourceList.js    ← Inventaire filtrable
    │   └── ResourceDetail.js  ← Fiche ressource + affectation
    ├── Suppliers/
    │   ├── SupplierList.js    ← Gestion des fournisseurs + blacklist
    │   └── SupplierRegister.js ← Inscription fournisseur (publique)
    ├── Breakdowns/
    │   ├── BreakdownList.js   ← Liste des tickets de panne
    │   └── BreakdownDetail.js ← Ticket + constat maintenance
    └── Departments/
        └── DepartmentList.js  ← Départements et membres
\end{verbatim}
\end{center}

\section{Gestion de l'État}

\subsection{AuthContext — État Global d'Authentification}

\begin{lstlisting}[style=js, caption={AuthContext.js — Gestion de l'état d'authentification}]
const AuthProvider = ({ children }) => {
    const [user, setUser] = useState(null);
    const [loading, setLoading] = useState(true);

    // Restauration de la session au demarrage
    useEffect(() => {
        const token = localStorage.getItem('token');
        const savedUser = localStorage.getItem('user');
        if (token && savedUser) {
            api.defaults.headers.common['Authorization'] =
                `Bearer ${token}`;
            setUser(JSON.parse(savedUser));
        }
        setLoading(false);
    }, []);

    const login = async (email, password) => {
        const { data } = await api.post('/auth/login',
            { email, password });
        // Persistance locale du token et du profil
        localStorage.setItem('token', data.token);
        localStorage.setItem('user', JSON.stringify(data.user));
        api.defaults.headers.common['Authorization'] =
            `Bearer ${data.token}`;
        setUser(data.user);
    };
    ...
};
\end{lstlisting}

\subsection{Intercepteur Axios}

L'intercepteur Axios centralise la gestion des erreurs d'authentification :

\begin{lstlisting}[style=js, caption={api.js — Intercepteur Axios}]
api.interceptors.response.use(
    response => response,  // Reponse valide: on passe directement
    error => {
        if (error.response?.status === 401) {
            // Token expire ou invalide: deconnexion automatique
            localStorage.removeItem('token');
            localStorage.removeItem('user');
            window.location.href = '/login';
        }
        return Promise.reject(error);
    }
);
\end{lstlisting}

\section{Routing et Protection des Routes}

Le composant \texttt{ProtectedRoute} protège toutes les pages authentifiées :

\begin{lstlisting}[style=js, caption={App.js — Protection des routes}]
const ProtectedRoute = ({ children }) => {
    const { user, loading } = useAuth();
    if (loading) return <div>Chargement...</div>;
    // Redirection vers /login si non authentifie
    return user ? children : <Navigate to="/login" />;
};
\end{lstlisting}

\section{Contrôle d'Accès Côté Interface}

Les éléments de l'interface sont conditionnés au rôle de l'utilisateur connecté, en complément du contrôle obligatoire côté serveur :

\begin{lstlisting}[style=js, caption={Exemple — Bouton conditionnel au rôle}]
// Le bouton "Nouveau" n'apparait que pour responsable et chef
{['responsable', 'chef_departement'].includes(user?.role) && (
    <Link to="/tenders/new" style={s.btn}>+ Nouveau</Link>
)}

// L'action "Blacklister" est reservee au responsable
{user?.role === 'responsable' && (
    <button onClick={() => handleBlacklist(supplier.id)}>
        Mettre en liste noire
    </button>
)}
\end{lstlisting}

%==========================================================
\chapter{Architecture de la Base de Données}
%==========================================================

\section{Schéma Relationnel}

La base de données PostgreSQL contient \textbf{11 tables} organisées autour des domaines métier du système.

\subsection{Tables Principales}

\begin{longtable}{|L{3.5cm}|L{5cm}|L{5cm}|}
\caption{Description des tables} \\
\hline
\rowcolor{primary}\textcolor{white}{\textbf{Table}} & \textcolor{white}{\textbf{Description}} & \textcolor{white}{\textbf{Colonnes clés}} \\
\hline
\endfirsthead
\hline
\textbf{Table} & \textbf{Description} & \textbf{Colonnes clés} \\
\hline
\endhead
users & Tous les utilisateurs du système & id (UUID PK), email (UNIQUE), password (hash), role (ENUM) \\
\hline
departments & Unités organisationnelles & id (UUID PK), name \\
\hline
suppliers & Profils des sociétés fournisseurs & id (UUID PK), user\_id (FK), company\_name, is\_blacklisted \\
\hline
tenders & Appels d'offre & id (UUID PK), start\_date, end\_date, status (ENUM), created\_by (FK) \\
\hline
tender\_needs & Besoins rattachés à un AO & tender\_id (FK), resource\_type, quantity, specs (JSONB) \\
\hline
offers & Propositions des fournisseurs & tender\_id (FK), supplier\_id (FK), total\_price, status (ENUM) \\
\hline
offer\_items & Détail article par article & offer\_id (FK), resource\_type, brand, unit\_price, quantity \\
\hline
resources & Inventaire matériel & id (UUID PK), inventory\_number (UNIQUE), status (ENUM), warranty\_end, specs (JSONB) \\
\hline
assignments & Affectations des ressources & resource\_id (FK), department\_id (FK), assigned\_user\_id (FK), is\_active \\
\hline
breakdowns & Tickets de panne & resource\_id (FK), reported\_by (FK), status (ENUM) \\
\hline
maintenance\_reports & Constats des techniciens & breakdown\_id (FK), technician\_id (FK), frequency (ENUM), order\_type (ENUM), can\_repair \\
\hline
\end{longtable}

\subsection{Extrait du Schéma SQL}

\begin{lstlisting}[style=sql, caption={Schéma — Tables principales (extrait)}]
-- Extension UUID
CREATE EXTENSION IF NOT EXISTS "uuid-ossp";

-- Utilisateurs avec roles
CREATE TABLE users (
    id            UUID PRIMARY KEY DEFAULT uuid_generate_v4(),
    name          VARCHAR(255) NOT NULL,
    email         VARCHAR(255) UNIQUE NOT NULL,
    password      VARCHAR(255) NOT NULL,
    role          VARCHAR(50) NOT NULL
                  CHECK (role IN (
                    'responsable','chef_departement',
                    'enseignant','fournisseur','technicien'
                  )),
    department_id UUID REFERENCES departments(id),
    created_at    TIMESTAMP DEFAULT NOW()
);

-- Appels d'offre avec cycle de vie
CREATE TABLE tenders (
    id          UUID PRIMARY KEY DEFAULT uuid_generate_v4(),
    title       VARCHAR(255) NOT NULL,
    description TEXT,
    start_date  DATE NOT NULL,
    end_date    DATE NOT NULL,
    status      VARCHAR(50) DEFAULT 'open'
                CHECK (status IN ('draft','open','closed','awarded')),
    created_by  UUID REFERENCES users(id),
    created_at  TIMESTAMP DEFAULT NOW()
);

-- Ressources avec specifications flexibles (JSONB)
CREATE TABLE resources (
    id               UUID PRIMARY KEY DEFAULT uuid_generate_v4(),
    inventory_number VARCHAR(100) UNIQUE NOT NULL,
    resource_type    VARCHAR(50) NOT NULL
                     CHECK (resource_type IN ('ordinateur','imprimante')),
    brand            VARCHAR(255),
    status           VARCHAR(50) DEFAULT 'available'
                     CHECK (status IN (
                       'available','assigned','maintenance','returned'
                     )),
    offer_id         UUID REFERENCES offers(id),
    supplier_id      UUID REFERENCES suppliers(id),
    warranty_end     DATE,
    specs            JSONB,
    created_at       TIMESTAMP DEFAULT NOW()
);
\end{lstlisting}

\section{Choix de Conception}

\subsection{Utilisation des UUID}

Toutes les clés primaires utilisent le type UUID (v4) plutôt que des entiers auto-incrémentés :
\begin{itemize}[leftmargin=1.5cm]
    \item \textbf{Sécurité} : les UUID ne sont pas prédictibles, réduisant les risques d'énumération
    \item \textbf{Distribution} : possibilité de générer des IDs côté client sans coordination avec la DB
    \item \textbf{Fusion de données} : facilite la migration ou la fusion de bases de données
\end{itemize}

\subsection{Utilisation du Type JSONB}

Les colonnes \texttt{specs} (ressources et besoins) utilisent le type JSONB de PostgreSQL :
\begin{itemize}[leftmargin=1.5cm]
    \item \textbf{Flexibilité} : les spécifications d'un ordinateur diffèrent de celles d'une imprimante sans nécessiter des tables séparées
    \item \textbf{Indexabilité} : JSONB peut être indexé avec GIN pour des recherches rapides
    \item \textbf{Contraintes} : les contraintes CHECK garantissent l'intégrité du reste des colonnes
\end{itemize}

\subsection{Gestion de l'Historique des Affectations}

La table \texttt{assignments} conserve l'historique complet grâce à la colonne \texttt{is\_active} :

\begin{lstlisting}[style=sql, caption={Requête — Affectation active d'une ressource}]
-- Affecter une ressource: desactiver l'affectation en cours
UPDATE assignments
SET is_active = FALSE, returned_date = NOW()
WHERE resource_id = $1 AND is_active = TRUE;

-- Creer la nouvelle affectation
INSERT INTO assignments
    (resource_id, department_id, assigned_user_id)
VALUES ($1, $2, $3);

-- Mettre a jour le statut de la ressource
UPDATE resources SET status = 'assigned' WHERE id = $1;
\end{lstlisting}

%==========================================================
\chapter{Architecture de Sécurité}
%==========================================================

\section{Vue d'Ensemble du Modèle de Sécurité}

Le système GRM implémente une sécurité en profondeur (\textit{defense in depth}) avec plusieurs couches de protection :

\begin{center}
\begin{verbatim}
  ┌──────────────────────────────────────────────┐
  │  COUCHE 1 — Transport (HTTPS/TLS)            │
  │  Chiffrement de toutes les communications    │
  └──────────────────────┬───────────────────────┘
                         ▼
  ┌──────────────────────────────────────────────┐
  │  COUCHE 2 — Authentification (JWT)           │
  │  Vérification de l'identité à chaque requête │
  └──────────────────────┬───────────────────────┘
                         ▼
  ┌──────────────────────────────────────────────┐
  │  COUCHE 3 — Autorisation (RBAC)              │
  │  Vérification du droit d'accès par rôle     │
  └──────────────────────┬───────────────────────┘
                         ▼
  ┌──────────────────────────────────────────────┐
  │  COUCHE 4 — Protection des Données           │
  │  Requêtes paramétrées (anti-injection SQL)   │
  │  Hachage bcrypt des mots de passe            │
  └──────────────────────────────────────────────┘
\end{verbatim}
\end{center}

\section{Authentification par Jeton JWT}

\subsection{Structure du Token}

Un token JWT se compose de trois parties encodées en Base64URL et séparées par des points :

\begin{center}
\texttt{header.payload.signature}
\end{center}

\textbf{Header :}
\begin{lstlisting}[style=js]
{ "alg": "HS256", "typ": "JWT" }
\end{lstlisting}

\textbf{Payload (claims) :}
\begin{lstlisting}[style=js]
{
  "id": "aaaaaaaa-aaaa-aaaa-aaaa-aaaaaaaaaaaa",
  "role": "responsable",
  "department_id": null,
  "supplier_id": null,
  "iat": 1715000000,   // Date d'emission (Unix timestamp)
  "exp": 1715086400    // Date d'expiration (24h plus tard)
}
\end{lstlisting}

\textbf{Signature :}
\begin{center}
\texttt{HMAC-SHA256(base64(header) + "." + base64(payload), JWT\_SECRET)}
\end{center}

\subsection{Middleware d'Authentification}

\begin{lstlisting}[style=js, caption={auth.js — Middleware authenticate}]
const authenticate = (req, res, next) => {
    const header = req.headers.authorization;

    // Verification de la presence du token
    if (!header)
        return res.status(401).json({ error: 'No token provided' });

    const token = header.split(' ')[1]; // Extraction du Bearer token
    try {
        // Verification de la signature et de l'expiration
        req.user = jwt.verify(token, process.env.JWT_SECRET);
        next(); // Transmission a la route suivante
    } catch {
        res.status(401).json({ error: 'Invalid token' });
    }
};
\end{lstlisting}

\section{Contrôle d'Accès par Rôle (RBAC)}

\subsection{Matrice des Autorisations}

\begin{longtable}{|L{4cm}|C{1.5cm}|C{1.5cm}|C{1.5cm}|C{1.5cm}|C{1.5cm}|}
\caption{Matrice RBAC complète} \\
\hline
\rowcolor{primary}\textcolor{white}{\textbf{Action}} & \textcolor{white}{\textbf{Resp.}} & \textcolor{white}{\textbf{Chef}} & \textcolor{white}{\textbf{Ensgt}} & \textcolor{white}{\textbf{Tech.}} & \textcolor{white}{\textbf{Fsseur}} \\
\hline
\endfirsthead
\hline
\textbf{Action} & \textbf{Resp.} & \textbf{Chef} & \textbf{Ensgt} & \textbf{Tech.} & \textbf{Fsseur} \\
\hline
\endhead
Créer un AO & ✓ & ✓ & & & \\
\hline
Fermer un AO & ✓ & & & & \\
\hline
Consulter les AO & ✓ & ✓ & & & ✓ \\
\hline
Soumettre une offre & & & & & ✓ \\
\hline
Accepter/Rejeter une offre & ✓ & & & & \\
\hline
Créer une ressource & ✓ & & & & \\
\hline
Modifier une ressource & ✓ & & & & \\
\hline
Supprimer une ressource & ✓ & & & & \\
\hline
Affecter une ressource & ✓ & & & & \\
\hline
Consulter les ressources & ✓ & ✓ & ✓ & & \\
\hline
Blacklister un fournisseur & ✓ & & & & \\
\hline
Consulter les fournisseurs & ✓ & & & & \\
\hline
Signaler une panne & & ✓ & ✓ & & \\
\hline
Rédiger un constat & & & & ✓ & \\
\hline
Décider retour fournisseur & ✓ & & & & \\
\hline
Résoudre une panne & ✓ & & & ✓ & \\
\hline
Gérer les départements & ✓ & & & & \\
\hline
\end{longtable}

\section{Protection contre les Injections SQL}

Toutes les interactions avec la base de données utilisent des requêtes paramétrées (\textit{prepared statements}) :

\begin{lstlisting}[style=js, caption={Bonne pratique — Requêtes paramétrées}]
// CORRECT: parametres separes de la requete
const { rows } = await pool.query(
    'SELECT * FROM users WHERE email = $1', [email]
);

// INCORRECT: concatenation dangereuse (injection SQL possible)
// await pool.query(`SELECT * FROM users WHERE email = '${email}'`);
\end{lstlisting}

%==========================================================
\chapter{Diagrammes de Séquence}
%==========================================================

\section{Diagramme 1 — Processus de Connexion}

\begin{verbatim}
Utilisateur    Navigateur (React)    Backend API         PostgreSQL
     │                │                   │                  │
     │──[Saisit email+mdp]──►             │                  │
     │                │──POST /api/auth/login─►              │
     │                │                   │──SELECT user WHERE email=$1──►│
     │                │                   │◄──{ user_row }───────────────│
     │                │                   │                  │
     │                │           [bcrypt.compare(mdp, hash)]│
     │                │                   │                  │
     │                │           [jwt.sign({id,role}, secret)]
     │                │                   │                  │
     │                │◄──200 { token, user }──              │
     │                │[localStorage.setItem(token)]         │
     │                │[Axios header: Bearer token]          │
     │◄──[Redirection Dashboard]          │                  │
\end{verbatim}

\section{Diagramme 2 — Soumission d'une Offre}

\begin{verbatim}
Fournisseur   React Frontend    Middleware JWT    offerController    PostgreSQL
     │               │                │                │                │
     │──[Formulaire]─►               │                │                │
     │               │──POST /api/offers/tender/:id──►│                │
     │               │               │──[Vérif JWT]   │                │
     │               │               │──[Rôle=fournisseur]            │
     │               │               │──────────────────►[connect()]───►│
     │               │               │                │──BEGIN          │
     │               │               │                │──SELECT is_blacklisted WHERE id=$1──►│
     │               │               │                │◄──{is_blacklisted: false}────────────│
     │               │               │                │──INSERT INTO offers VALUES(...)──────►│
     │               │               │                │◄──{ offer_row }──────────────────────│
     │               │               │                │──INSERT INTO offer_items (loop)──────►│
     │               │               │                │──COMMIT         │
     │               │               │                │──[release()]    │
     │               │◄──201 { offer }─────────────────               │
     │◄──[Confirmation]              │                │                │
\end{verbatim}

\section{Diagramme 3 — Signalement et Traitement d'une Panne}

\begin{verbatim}
Enseignant    Technicien    Responsable    Backend API         PostgreSQL
     │             │              │              │                  │
     │──[Signale panne]─────────────────────────►                  │
     │                            │    POST /api/breakdowns         │
     │                            │              │──INSERT breakdowns──►│
     │                            │              │──UPDATE resources SET status='maintenance'─►│
     │                            │              │◄──201 { breakdown }──────────────────────│
     │                            │              │                  │
     │          [Technicien consulte pannes]      │                  │
     │             │──GET /api/breakdowns─────────►                 │
     │             │◄──[liste pannes ouvertes]─────                 │
     │             │                              │                  │
     │             │──[Rédige constat]──────────────────────────────►│
     │             │  POST /api/breakdowns/:id/maintenance-report   │
     │             │◄──201 { report }──────────────────────────────│
     │             │                              │                  │
     │                            │──[Consulte constat]──────────────►│
     │                            │  GET /api/breakdowns/:id        │
     │                            │◄──{ breakdown + reports }────────│
     │                            │                                  │
     │                            │──[Décide retour fournisseur]──────►│
     │                            │  PATCH /api/breakdowns/:id/return-to-supplier
     │                            │              │──UPDATE breakdowns SET status='returned_...'─►│
     │                            │              │──UPDATE resources SET status='returned'──────►│
     │                            │◄──200 { message }──────────────│
\end{verbatim}

%==========================================================
\chapter{Architecture de Déploiement}
%==========================================================

\section{Environnements}

\begin{longtable}{|L{2.5cm}|L{3.5cm}|L{3.5cm}|L{4cm}|}
\caption{Description des environnements} \\
\hline
\rowcolor{primary}\textcolor{white}{\textbf{Env.}} & \textcolor{white}{\textbf{Usage}} & \textcolor{white}{\textbf{Configuration}} & \textcolor{white}{\textbf{Base de données}} \\
\hline
\endfirsthead
\hline
\textbf{Env.} & \textbf{Usage} & \textbf{Configuration} & \textbf{Base de données} \\
\hline
\endhead
Développement & Développement local & \texttt{.env} local, hot-reload (nodemon + CRA) & PostgreSQL local / Docker \\
\hline
Test & Tests automatisés & Variables d'env. dédiées & Base de données de test isolée \\
\hline
Production & Déploiement final & Docker Compose, HTTPS & PostgreSQL avec volume persistant \\
\hline
\end{longtable}

\section{Diagramme de Déploiement Docker}

\begin{center}
\begin{verbatim}
  ╔══════════════════════════════════════════════════════════╗
  ║              Hôte Docker (Serveur Linux)                 ║
  ║                                                          ║
  ║  ┌──────────────────────────────────────────────────┐   ║
  ║  │           docker-compose.yml                     │   ║
  ║  │                                                  │   ║
  ║  │  ┌─────────────────┐    ┌──────────────────┐    │   ║
  ║  │  │  service: db    │    │ service: backend │    │   ║
  ║  │  │  postgres:15    │◄───│ node:20-alpine   │    │   ║
  ║  │  │  port: 5432     │    │ port: 5000       │    │   ║
  ║  │  │  volume: pgdata │    │ ENV: DB_HOST=db  │    │   ║
  ║  │  └─────────────────┘    └────────▲─────────┘    │   ║
  ║  │                                  │               │   ║
  ║  │  ┌───────────────────────────────┴───────────┐  │   ║
  ║  │  │  service: frontend                        │  │   ║
  ║  │  │  nginx:alpine — port: 80                  │  │   ║
  ║  │  │  /usr/share/nginx/html ← build React      │  │   ║
  ║  │  │  location /api → proxy_pass backend:5000  │  │   ║
  ║  │  └───────────────────────────────────────────┘  │   ║
  ║  └──────────────────────────────────────────────────┘   ║
  ║                                                          ║
  ╚══════════════════════════════════════════════════════════╝
                          │ Port 80 (HTTP) / 443 (HTTPS)
                    [Internet / Réseau Faculté]
                          │
               [Navigateur Client]
\end{verbatim}
\end{center}

\section{Configuration Docker Compose}

\begin{lstlisting}[style=bash, caption={docker-compose.yml — Configuration complète}]
version: '3.8'
services:
  # Service PostgreSQL avec persistance
  db:
    image: postgres:15-alpine
    restart: unless-stopped
    environment:
      POSTGRES_DB: grm_db
      POSTGRES_USER: postgres
      POSTGRES_PASSWORD: ${DB_PASSWORD}
    volumes:
      - ./backend/database/schema.sql:/docker-entrypoint-initdb.d/schema.sql
      - pgdata:/var/lib/postgresql/data
    healthcheck:
      test: ["CMD-SHELL", "pg_isready -U postgres"]
      interval: 5s
      retries: 5

  # Service API Node.js
  backend:
    build: ./backend
    restart: unless-stopped
    ports: ["5000:5000"]
    environment:
      DB_HOST: db           # Nom du service Docker
      JWT_SECRET: ${JWT_SECRET}
    depends_on:
      db:
        condition: service_healthy  # Attend que PostgreSQL soit pret

  # Service Frontend (Nginx + build React)
  frontend:
    build: ./frontend
    restart: unless-stopped
    ports: ["80:80"]
    depends_on: [backend]

volumes:
  pgdata:   # Volume nomme pour la persistance des donnees
\end{lstlisting}

\section{Build Multi-Étapes du Frontend}

\begin{lstlisting}[style=bash, caption={frontend/Dockerfile — Build multi-étapes}]
# Etape 1: Build de l'application React
FROM node:20-alpine AS build
WORKDIR /app
COPY package*.json ./
RUN npm ci                   # Installation reproductible (lockfile)
COPY . .
RUN npm run build            # Production build optimise

# Etape 2: Serveur Nginx pour le build statique
FROM nginx:alpine
# Copie uniquement les artefacts de build (pas les sources)
COPY --from=build /app/build /usr/share/nginx/html
COPY nginx.conf /etc/nginx/conf.d/default.conf
EXPOSE 80
\end{lstlisting}

\section{Configuration Nginx}

\begin{lstlisting}[style=bash, caption={nginx.conf — Configuration du reverse proxy}]
server {
    listen 80;
    root /usr/share/nginx/html;
    index index.html;

    # SPA: toutes les routes non-trouvees renvoient index.html
    # (React Router gere le routing cote client)
    location / {
        try_files $uri $uri/ /index.html;
    }

    # Proxy des appels API vers le backend Node.js
    location /api {
        proxy_pass http://backend:5000;
        proxy_http_version 1.1;
        proxy_set_header Host $host;
        proxy_set_header X-Real-IP $remote_addr;
        proxy_set_header X-Forwarded-For $proxy_add_x_forwarded_for;
    }
}
\end{lstlisting}

%==========================================================
\chapter{Stratégie de Tests}
%==========================================================

\section{Niveaux de Tests}

\begin{longtable}{|L{3cm}|L{4.5cm}|L{4.5cm}|L{2cm}|}
\caption{Stratégie de tests} \\
\hline
\rowcolor{primary}\textcolor{white}{\textbf{Niveau}} & \textcolor{white}{\textbf{Description}} & \textcolor{white}{\textbf{Outils}} & \textcolor{white}{\textbf{Couverture}} \\
\hline
\endfirsthead
\hline
\textbf{Niveau} & \textbf{Description} & \textbf{Outils} & \textbf{Couverture} \\
\hline
\endhead
Tests Unitaires & Test de fonctions isolées (middleware auth, validations) & Jest, jest-mock & Middleware, utilitaires \\
\hline
Tests d'Intégration & Test des endpoints API avec une vraie DB de test & Supertest + Jest & Tous les endpoints \\
\hline
Tests E2E & Simulation du parcours utilisateur complet & Cypress & Scénarios critiques \\
\hline
Tests Manuels & Validation fonctionnelle avec les comptes de test & Navigateur & Interface utilisateur \\
\hline
\end{longtable}

\section{Exemples de Tests d'Intégration}

\begin{lstlisting}[style=js, caption={Tests d'intégration — Authentification}]
describe('POST /api/auth/login', () => {
    it('retourne un JWT valide pour des credentials corrects', async () => {
        const res = await request(app)
            .post('/api/auth/login')
            .send({
                email: 'responsable@faculte.ma',
                password: 'admin123'
            });

        expect(res.status).toBe(200);
        expect(res.body).toHaveProperty('token');
        expect(res.body.user.role).toBe('responsable');
    });

    it('retourne 401 pour un mot de passe incorrect', async () => {
        const res = await request(app)
            .post('/api/auth/login')
            .send({ email: 'responsable@faculte.ma', password: 'wrong' });

        expect(res.status).toBe(401);
        expect(res.body).toHaveProperty('error');
    });
});

describe('POST /api/tenders', () => {
    it('refuse la creation sans authentification', async () => {
        const res = await request(app)
            .post('/api/tenders')
            .send({ title: 'Test AO' });

        expect(res.status).toBe(401);
    });

    it('refuse la creation si role != responsable/chef', async () => {
        // Authentification en tant que fournisseur
        const loginRes = await request(app)
            .post('/api/auth/login')
            .send({ email: 'fournisseur@test.ma', password: 'mdp' });

        const res = await request(app)
            .post('/api/tenders')
            .set('Authorization', `Bearer ${loginRes.body.token}`)
            .send({ title: 'AO non autorise' });

        expect(res.status).toBe(403);
    });
});
\end{lstlisting}

%==========================================================
\chapter{Guide de Démarrage}
%==========================================================

\section{Prérequis}

\begin{longtable}{|L{4cm}|L{4cm}|L{6cm}|}
\caption{Prérequis d'installation} \\
\hline
\rowcolor{primary}\textcolor{white}{\textbf{Outil}} & \textcolor{white}{\textbf{Version minimale}} & \textcolor{white}{\textbf{Usage}} \\
\hline
\endfirsthead
\hline
\textbf{Outil} & \textbf{Version minimale} & \textbf{Usage} \\
\hline
\endhead
Node.js & 18.x LTS & Exécution du backend et build frontend \\
\hline
npm & 9.x & Gestionnaire de paquets JavaScript \\
\hline
PostgreSQL & 14.x & Base de données relationnelle \\
\hline
Docker & 24.x & Conteneurisation (optionnel en dev) \\
\hline
Docker Compose & 2.x & Orchestration multi-conteneurs \\
\hline
Git & 2.x & Versionnement du code source \\
\hline
\end{tabularx}

\section{Installation en Mode Développement}

\begin{lstlisting}[style=bash, caption={Installation manuelle — Développement}]
# 1. Cloner le depot
git clone https://github.com/grm-project/grm.git
cd grm

# 2. Initialiser la base de donnees PostgreSQL
psql -U postgres -c "CREATE DATABASE grm_db;"
psql -U postgres -d grm_db -f backend/database/schema.sql

# 3. Configurer le backend
cd backend
cp .env.example .env
# Editer .env: DB_PASSWORD, JWT_SECRET

npm install
npm run dev
# Serveur API: http://localhost:5000

# 4. Configurer le frontend (dans un nouveau terminal)
cd ../frontend
npm install
npm start
# Application: http://localhost:3000
\end{lstlisting}

\section{Déploiement en Production (Docker)}

\begin{lstlisting}[style=bash, caption={Déploiement production avec Docker Compose}]
# 1. Configurer les variables de production
cp .env.example .env.prod
# Editer: DB_PASSWORD fort, JWT_SECRET aleatoire 256 bits

# 2. Lancer tous les services
docker-compose --env-file .env.prod up --build -d

# 3. Verification du statut
docker-compose ps
# db         running (healthy)
# backend    running
# frontend   running

# 4. Application disponible
# http://votre-serveur.ma
\end{lstlisting}

\section{Comptes de Test Préconfigurés}

\begin{longtable}{|L{5cm}|L{3cm}|L{3cm}|L{3cm}|}
\caption{Comptes de test disponibles après initialisation} \\
\hline
\rowcolor{primary}\textcolor{white}{\textbf{Email}} & \textcolor{white}{\textbf{Mot de passe}} & \textcolor{white}{\textbf{Rôle}} & \textcolor{white}{\textbf{Département}} \\
\hline
\endfirsthead
\hline
\textbf{Email} & \textbf{Mot de passe} & \textbf{Rôle} & \textbf{Département} \\
\hline
\endhead
responsable@faculte.ma & admin123 & responsable & — \\
\hline
chef.info@faculte.ma & admin123 & chef\_departement & Informatique \\
\hline
ahmed.benali@faculte.ma & admin123 & enseignant & Informatique \\
\hline
tech@faculte.ma & admin123 & technicien & — \\
\hline
\end{longtable}

\end{document}
